\section{Discussion}
The observation protocol is part of a predictive system, not merely a fixed
preprocessing choice. When the target summarises a longer trajectory than the
measurements supplied to the learner, performance reflects both the learner
and the information admitted by the protocol. More training units and more
flexible predictors can reduce error under the realised protocol, but they do
not reveal the value of measurements that were never collected. Even the
population measurement--target law can leave dependence unresolved that is
decisive for an alternative protocol.

The practical meaning of ``more data'' therefore depends on how those data are
collected. If the ambiguity is invisible to the realised protocol, enlarging
the sample under that protocol reproduces the same constraints. A smaller
sample observed more extensively may be more useful. Targeted augmentation
removes ambiguity affecting one specified alternative, whereas dense
calibration supports a broader family. This motivates a large routine cohort
accompanied by a smaller, intensively observed calibration subset chosen with
future protocol decisions in mind.

Finite calibration determines how finely protocols can be compared. The
precision of an optimiser's output should not be confused with the precision
supported by the data. \Cref{fig:calibration} shows that richer candidate
classes can reduce selection regret before all candidate values are estimated
accurately, because selection depends mainly on ordering protocols near the
optimum. Useful granularity is set by the value gaps among competitive
protocols and the information available to resolve them.

The retrospective analyses illustrate this distinction. In Sleep, broad
dispersed-versus-contiguous contrasts are more reproducible than exact learned
supports, whose locations and held-out advantages vary across targets, source
studies and subsamples (\cref{fig:real,fig:sweep}). In AF, distributing a
matched budget across the record is substantially more predictive of
full-record burden than concentrating it in one block at the multi-window
budgets considered. This is not a general rule in favour of dispersion; equal
observation time need not carry equal information for a target defined over a
longer horizon. Protocol quality is also target-dependent: \cref{tab:s5}
shows that observing the process well and observing it well for a specified
aggregate can lead to different designs.

The same problem arises beyond temporal sampling. Short ECG recordings may not
reveal the cross-day dependence needed to evaluate dispersed monitoring; an
overlapping multimodal subset may be needed before imaging, proteomic or
wearable panels can be valued; and current environmental sensors may not reveal
the spatial dependence governing the value of new locations. These are not
ordinary feature-selection problems, because the proposed measurements may
never have been observed jointly with the target. Calibration precedes the
usual optimisation problem in longitudinal and sensor design
\citep{chaloner1995bayesian,ji2017optimal,krause2008near,joshi2009sensor}.

When point identification fails, the range of compatible values can still
support partial comparisons and robust design based on worst-case value or
minimax regret \citep{manski2003partial,tamer2010partial}. The Gaussian analysis
exposes this ambiguity through covariance directions, but the principle is
broader: an alternative is evaluable whenever its value is constant over the
latent laws compatible with the realised data. In more general models, the
corresponding task is to characterise that compatible set and collect
measurements that shrink the decision-relevant value range.

\section{Conclusion}

Counterfactual protocol evaluation determines whether data from one observation
system fix the predictive value of another that was never deployed. Existing
data need not determine the value of measurements that were never made;
target-relevant dependence may remain observationally unconstrained even with
an infinite benchmark. Value-specific identification isolates the ambiguity
that matters, targeted augmentation can remove it without recovering the full
latent system, finite calibration determines which value gaps are resolvable,
and exact marginal gains support cost-constrained design at that resolution.
Identification should therefore precede optimisation, both for temporal
protocols and for observation systems more generally.